\section{Related Work}
In this section, we reflect on the prior literature in human-AI interaction and situate, within its canon, the agenda of building believable proxies of human behavior. This agenda, once hailed as a north star in the interaction, game, and artificial intelligence communities~\cite{4_laird2001human, 6_riedl2012interactive, 21_riedl2005objective, 22_bates1994role}, has remained challenging due to the complexity of human behavior~\cite{23_brooks2000cog, 25_yannakakis2012game}. We synthesize this research to suggest that large language models, though not sufficient by themselves, open up a new angle for creating believable agents when leveraged using the appropriate architecture.

\subsection{Human-AI Interaction}
Interactive artificial intelligence systems aim to combine human insights and capabilities in computational artifacts that can augment their users ~\cite{amershi2014power, fails2003interactive}. A long line of work has explored ways to enable users to interactively specify model behavior. For instance, Crayons demonstrated an early vision of interactive machine learning, allowing non-expert users to train classifiers~\cite{fails2003interactive}. Further work helped to articulate how end users might describe their classification goals to the system through examples~\cite{Fogarty2008} or demonstration~\cite{fiebrink2010wekinator}. Recent advancements have extended these explorations to deep learning~\cite{lam2023model} and prompt-based authoring~\cite{jiang_PromptMaker,62_Wu2021AIChains,liu2022opal}.

Meanwhile, a persistent thread of research has advanced the case for language- and agent-based interaction in human-computer interaction. Formative work such as SHRDLU~\cite{winograd1971procedures} and ELIZA~\cite{weizenbaum1966eliza} demonstrated the opportunities and the risks associated with natural language interaction with computing systems. As research progressed, it became evident that autonomous agents could offer new metaphors for delegation and interaction~\cite{MaesPattie1995}, but the boundaries of delegation between humans and agents have remained the subject of ongoing debate and refinement~\cite{shneiderman1997direct,horvitz1999principles,shneiderman2022human}. Recently, this technology has reached a level of stability that enables agents to interact via natural language in large and complex online social environments (e.g.,~\cite{Krishna2022}). Natural language interaction provides a novel modality that can enhance user abilities in domains such as photo editing~\cite{laput2013pixeltone, fourney2011query, adar2014command} and code editing~\cite{rong2016codemend}.

We convene these threads of work to show that we can now create agents that proxy human behavior for interactive systems, and interact with them using natural language. In doing so, this work reopens the door to examining foundational human-computer interaction questions around cognitive models such as GOMS and Keystroke-Level Model (KLM)~\cite{card1983psychology,card1980keystroke}, around prototyping tools~\cite{9_park2022socialsimulacra}, and around ubiquitous computing applications~\cite{weiser1991computer,dey2001understanding,fast2016augur}.

\subsection{Believable Proxies of Human Behavior}
Prior literature has described \textit{believability}, or \textit{believable agents}, as a central design and engineering goal. Believable agents are designed to provide an illusion of life and present a facade of realism in the way they appear to make decisions and act on their own volition, similar to the characters in Disney movies~\cite{22_bates1994role, 42_Thomas1981Disney}. These agents can populate and perceive an open world environment like the one we inhabit~\cite{4_laird2001human, 22_bates1994role}, and strive to behave in ways that exhibit emergent behaviors grounded in social interactions with users or other agents with the aim of becoming believable proxies of our behavior in hypothetical simulations of individuals and communities~\cite{26_mccoy2012prom, 43_Burkinshaw2009Alice, 44_Francis2010Minecraft}. Historically, these agents were developed in the context of intelligent game non-player characters (NPCs)~\cite{4_laird2001human, 6_riedl2012interactive}. Creating NPCs with believable behavior, if possible, could enhance player experiences in games and interactive fictions by enabling emergent narratives~\cite{1_swartout2006toward, 2_aylett1999narrative, 3_brenner2010creating, ibister2000consistency} and social interactions with the agents~\cite{5_zubek2002towards}. However, more importantly, game worlds provide increasingly realistic representations of real-world affordances, and as observed by Laird and van Lent in 2001, these simulated worlds offer accessible testbeds for developers of believable agents to finesse the agents' cognitive capabilities without worrying about implementing robotics in the real world or creating simulation environments from scratch~\cite{4_laird2001human, 6_riedl2012interactive}.
 
A diverse set of approaches to creating believable agents emerged over the past four decades.
In implementation, however, these approaches often simplified the environment or dimensions of agent behavior to make the effort more manageable~\cite{23_brooks2000cog, 39_Minsky1970Draft}. Rule-based approaches, such as finite-state machines~\cite{30_NEURIPS2021_86e8f7ab, 32_Umarov2012Believable} and behavior trees~\cite{45_Knafla2011Introduction, 46_Pillosu2009Coordinating, 47_Hecker2011Spore} account for the brute force approach of human-authoring the agent's behavior~\cite{26_mccoy2012prom}. They provide a straightforward way of creating simple agents that is still the most dominant approach today~\cite{37_McCoy2009Comme, 24_miyashita2017developing, 25_yannakakis2012game}, and can even handle rudimentary social interactions, as shown in games such as Mass Effect~\cite{48_BioWare2007Mass} and The Sims~\cite{49_ElectronicArts2009Sims} series. Nonetheless, manually crafting behavior that can comprehensively address the breadth of possible interactions in an open world is untenable. This means that the resulting agent behaviors may not fully represent the consequences of their interactions~\cite{26_mccoy2012prom, 27_mccoy2011prom, 28_mccoy2011comme}, and cannot perform new procedures that were not hard-coded in their script~\cite{30_NEURIPS2021_86e8f7ab, 32_Umarov2012Believable}. On the other hand, prevalent learning-based approaches for creating believable agents, such as reinforcement learning, have overcome the challenge of manual authoring by letting the agents learn their behavior, and have achieved superhuman performance in recent years in games such as AlphaStar for Starcraft~\cite{50_Vinyals2019Grandmaster} and OpenAI Five for Dota 2~\cite{51_Berner2019Dota}. However, their success has largely taken place in adversarial games with readily definable rewards that a learning algorithm can optimize for. They have not yet addressed the challenge of creating believable agents in an open world~\cite{30_NEURIPS2021_86e8f7ab, 24_miyashita2017developing, 38_Hausknecht2020Interactive}.

Cognitive architectures in computation, pioneered by Newell, aimed to build the infrastructure for supporting a comprehensive set of cognitive functions~\cite{14_newell1990unified} that suited the all-encompassing nature of believable agents held in its original vision. They fueled some of the earliest examples of believable agents. For instance, Quakebot-SOAR~\cite{54_Laird2000It} and ICARUS~\cite{33_LangleyIcarus, 34_Choi2021Believable} generated NPCs in first-person shooter games, while TacAir-SOAR~\cite{55_Pew1998Modeling} generated pilots in aerial combat training simulations. The architectures used by these agents differed (Quakebot- and TacAir-SOAR relied on SOAR~\cite{56_Laird2012Soar}, while ICARUS relied on its own variation that was inspired by SOAR and ACT-R~\cite{57_Anderson1993Rules}), but they shared the same underlying principle~\cite{58_Laird2017Standard}. They maintained short-term and long-term memories, filled these memories with symbolic structures, and operated in perceive-plan-act cycles, dynamically perceiving the environment and matching it with one of the manually crafted action procedures~\cite{32_Umarov2012Believable, 36_Laird2001Knows}. Agents created using cognitive architectures aimed to be generalizable to most, if not all, open world contexts and exhibited robust behavior for their time. However, their space of action was limited to manually crafted procedural knowledge, and they did not offer a mechanism through which the agents could be inspired to seek new behavior. As such, these agents were deployed mostly in non-open world contexts such as first-person shooter games~\cite{54_Laird2000It, 34_Choi2021Believable} or blocks worlds~\cite{33_LangleyIcarus}.

Today, creating believable agents as described in its original definition remains an open problem~\cite{25_yannakakis2012game, 6_riedl2012interactive}. Many have moved on, arguing that although current approaches for creating believable agents might be cumbersome and limited, they are good enough to support existing gameplay and interactions~\cite{25_yannakakis2012game, 52_Champandard2012Tutorial, 53_Nareyek2007Game}. Our argument is that large language models offer an opportunity to re-examine these questions, provided that we can craft an effective architecture to synthesize memories into believable behavior. We offer a step toward such an architecture in this paper. 


\subsection{Large Language Models and Human Behavior}
Generative agents leverage a large language model to power their behavior. The key observation is that large language models encode a wide range of human behavior from their training data~\cite{59_brown2020language, 60_bommasani2022opportunities}. If prompted with a narrowly defined context, the models can be used to generate believable behavior. Recent work has demonstrated the efficacy of this approach. For instance, social simulacra used a large language model to generate users that would populate new social computing systems to prototype their emergent social dynamics~\cite{9_park2022socialsimulacra}. This approach used a prompt chain~\cite{61_Wu2022PromptChainer, 62_Wu2021AIChains} to generate short natural language descriptions of personas and their behaviors as they appear in the system being prototyped. Other empirical studies have replicated existing social science studies~\cite{63_horton2023large}, political surveys~\cite{64_Sorensen_2022}, and generated synthetic data~\cite{hamalainen2023evaluating}. Large language models have also been used to generate interactive human behavior for users to engage with. In gaming, for instance, these models have been employed to create interactive fiction~\cite{65_Freiknecht2020Procedural} and text adventure games~\cite{66_callison-burch-etal-2022-dungeons}. With their ability to generate and decompose action sequences, large language models have also been used in planning robotics tasks~\cite{67_huang2022inner}. For example, when presented with a task, such as picking up a bottle, the model is prompted to break down the task into smaller action sequences, such as heading to the table where the bottle is located and picking it up.

We posit that, based on the work summarized above, large language models can become a key ingredient for creating believable agents. The existing literature largely relies on what could be considered first-order templates that employ few-shot prompts~\cite{68_Gao2020Making, 69_Liu2021What} or chain-of-thought prompts~\cite{70_wei2023chainofthought}. These templates are effective in generating behavior that is conditioned solely on the agent's current environment (e.g., how would a troll respond to a given post, what actions would a robot need to take to enter a room given that there is a door). However, believable agents require conditioning not only on their current environment but also on a vast amount of past experience, which is a poor fit (and as of today, impossible due to the underlying models' limited context window) using first-order prompting. Recent studies have attempted to go beyond first-order prompting by augmenting language models with a static knowledge base and an information retrieval scheme~\cite{khattab2023demonstratesearchpredict} or with a simple summarization scheme~\cite{wu2021recursively}. This paper extends these ideas to craft an agent architecture that handles retrieval where past experience is dynamically updated at each time step and mixed with agents' current context and plans, which may either reinforce or contradict each other.